
\documentclass{article}
\usepackage{colortbl}
\usepackage{makecell}
\usepackage{multirow}
\usepackage{supertabular}

\begin{document}

\newcounter{utterance}

\centering \large Interaction Transcript for game `matchit\_ascii', experiment `similar\_grid\_2', episode 2 with gemma{-}4{-}26b{-}a4b{-}it.
\vspace{24pt}

{ \footnotesize  \setcounter{utterance}{1}
\setlength{\tabcolsep}{0pt}
\begin{supertabular}{c@{$\;$}|p{.15\linewidth}@{}p{.15\linewidth}p{.15\linewidth}p{.15\linewidth}p{.15\linewidth}p{.15\linewidth}}
    \# & $\;$A & \multicolumn{4}{c}{Game Master} & $\;\:$B\\
    \hline

    \theutterance \stepcounter{utterance}  
    & & \multicolumn{4}{p{0.6\linewidth}}{
        \cellcolor[rgb]{0.9,0.9,0.9}{
            \makecell[{{p{\linewidth}}}]{
                \texttt{\tiny{[P1$\langle$GM]}}
                \texttt{Egy együttműködésen alapuló találgatós játékban veszel részt.} \\
\texttt{A cél annak kiderítése, hogy ez a rács és egy másik kép, amelyet csak én látok, ugyanaz{-}e.} \\
\texttt{Kérlek, először írd le a te rácsodat. Ezután én megadom az én leírásomat, és kérdéseket tehetünk fel egymásnak a rácsokról, hogy kiderítsük, ugyanazok{-}e.  Most kezdd a rövid rácsleírásodat "LEÍRÁS:" jelzővel, majd add meg a leírást. Ne adj hozzá semmi mást. Ez a te rácsod:} \\
\texttt{X ▢ ▢ ▢ X} \\
\texttt{X X ▢ X X} \\
\texttt{X ▢ X ▢ X} \\
\texttt{X ▢ ▢ ▢ X} \\
\texttt{X ▢ ▢ ▢ X} \\
            }
        }
    }
    & \\ \\

    \theutterance \stepcounter{utterance}  
    & \multicolumn{4}{p{0.6\linewidth}}{
        \cellcolor[rgb]{0.9,0.9,0.9}{
            \makecell[{{p{\linewidth}}}]{
                \texttt{\tiny{[P1$\rangle$GM]}}
                \texttt{LEÍRÁS:} \\
\texttt{X ▢ ▢ ▢ X} \\
\texttt{X X ▢ X X} \\
\texttt{X ▢ X ▢ X} \\
\texttt{X ▢ ▢ ▢ X} \\
\texttt{X ▢ ▢ ▢ X} \\
            }
        }
    }
    & & \\ \\

    \theutterance \stepcounter{utterance}  
    & & & \multicolumn{2}{p{0.3\linewidth}}{
        \cellcolor[rgb]{0.9,0.9,0.9}{
            \makecell[{{p{\linewidth}}}]{
                \texttt{\tiny{[GM$|$GM]}}
                \texttt{continue} \\
            }
        }
    }
    & & \\ \\

    \theutterance \stepcounter{utterance}  
    & & \multicolumn{4}{p{0.6\linewidth}}{
        \cellcolor[rgb]{0.9,0.9,0.9}{
            \makecell[{{p{\linewidth}}}]{
                \texttt{\tiny{[GM$\rangle$P2]}}
                \texttt{Egy együttműködésen alapuló találgatós játékban veszel részt.} \\
\texttt{A cél annak kiderítése, hogy ez a rács és egy másik kép, amelyet csak én látok, ugyanaz{-}e.} \\
\texttt{Kérlek, először írd le a te rácsodat. Ezután én megadom az én leírásomat, és kérdéseket tehetünk fel egymásnak a rácsokról, hogy kiderítsük, ugyanazok{-}e.  Most kezdd a rövid rácsleírásodat "LEÍRÁS:" jelzővel, majd add meg a leírást. Ne adj hozzá semmi mást. Ez a te rácsod:} \\
\texttt{X ▢ ▢ ▢ X} \\
\texttt{X ▢ ▢ X X} \\
\texttt{X X X ▢ X} \\
\texttt{X ▢ ▢ ▢ X} \\
\texttt{X ▢ ▢ ▢ X} \\
            }
        }
    }
    & \\ \\

    \theutterance \stepcounter{utterance}  
    & & & \multicolumn{4}{p{0.6\linewidth}}{
        \cellcolor[rgb]{0.9,0.9,0.9}{
            \makecell[{{p{\linewidth}}}]{
                \texttt{\tiny{[GM$\langle$P2]}}
                \texttt{LEÍRÁS:} \\
\texttt{X ▢ ▢ ▢ X} \\
\texttt{X ▢ ▢ X X} \\
\texttt{X X X ▢ X} \\
\texttt{X ▢ ▢ ▢ X} \\
\texttt{X ▢ ▢ ▢ X} \\
            }
        }
    }
     \\ \\

    \theutterance \stepcounter{utterance}  
    & & & \multicolumn{2}{p{0.3\linewidth}}{
        \cellcolor[rgb]{0.9,0.9,0.9}{
            \makecell[{{p{\linewidth}}}]{
                \texttt{\tiny{[GM$|$GM]}}
                \texttt{continue} \\
            }
        }
    }
    & & \\ \\

    \theutterance \stepcounter{utterance}  
    & & \multicolumn{4}{p{0.6\linewidth}}{
        \cellcolor[rgb]{0.9,0.9,0.9}{
            \makecell[{{p{\linewidth}}}]{
                \texttt{\tiny{[GM$\rangle$P2]}}
                \texttt{Ez az énLEÍRÁS:} \\
\texttt{X ▢ ▢ ▢ X} \\
\texttt{X X ▢ X X} \\
\texttt{X ▢ X ▢ X} \\
\texttt{X ▢ ▢ ▢ X} \\
\texttt{X ▢ ▢ ▢ X} \\
\texttt{Most tegyél fel egy kérdést annak érdekében, hogy új aspektusokat deríts ki az én rácsomról, amely eltérhet a te rácsodtól. Kezdd "KÉRDÉS:" jelzővel, és ne adj hozzá semmi mást.} \\
            }
        }
    }
    & \\ \\

    \theutterance \stepcounter{utterance}  
    & & & \multicolumn{4}{p{0.6\linewidth}}{
        \cellcolor[rgb]{0.9,0.9,0.9}{
            \makecell[{{p{\linewidth}}}]{
                \texttt{\tiny{[GM$\langle$P2]}}
                \texttt{KÉRDÉS:} \\
\texttt{A második sorban az X karakterek elhelyezkedése pontosan milyen sorrendben van?} \\
            }
        }
    }
     \\ \\

    \theutterance \stepcounter{utterance}  
    & & & \multicolumn{2}{p{0.3\linewidth}}{
        \cellcolor[rgb]{0.9,0.9,0.9}{
            \makecell[{{p{\linewidth}}}]{
                \texttt{\tiny{[GM$|$GM]}}
                \texttt{continue} \\
            }
        }
    }
    & & \\ \\

    \theutterance \stepcounter{utterance}  
    & & \multicolumn{4}{p{0.6\linewidth}}{
        \cellcolor[rgb]{0.9,0.9,0.9}{
            \makecell[{{p{\linewidth}}}]{
                \texttt{\tiny{[P1$\langle$GM]}}
                \texttt{Ez az énLEÍRÁS:} \\
\texttt{X ▢ ▢ ▢ X} \\
\texttt{X ▢ ▢ X X} \\
\texttt{X X X ▢ X} \\
\texttt{X ▢ ▢ ▢ X} \\
\texttt{X ▢ ▢ ▢ X} \\
\texttt{KÉRDÉS:} \\
\texttt{A második sorban az X karakterek elhelyezkedése pontosan milyen sorrendben van?Kezdd a válaszodat a "VÁLASZ:" jelzővel, és ne adj hozzá semmi mást.} \\
            }
        }
    }
    & \\ \\

    \theutterance \stepcounter{utterance}  
    & \multicolumn{4}{p{0.6\linewidth}}{
        \cellcolor[rgb]{0.9,0.9,0.9}{
            \makecell[{{p{\linewidth}}}]{
                \texttt{\tiny{[P1$\rangle$GM]}}
                \texttt{VÁLASZ:} \\
\texttt{Az én második soromban az X karakterek az első, második, negyedik és ötödik pozícióban vannak (X X ▢ X X).} \\
            }
        }
    }
    & & \\ \\

    \theutterance \stepcounter{utterance}  
    & & & \multicolumn{2}{p{0.3\linewidth}}{
        \cellcolor[rgb]{0.9,0.9,0.9}{
            \makecell[{{p{\linewidth}}}]{
                \texttt{\tiny{[GM$|$GM]}}
                \texttt{continue} \\
            }
        }
    }
    & & \\ \\

    \theutterance \stepcounter{utterance}  
    & & \multicolumn{4}{p{0.6\linewidth}}{
        \cellcolor[rgb]{0.9,0.9,0.9}{
            \makecell[{{p{\linewidth}}}]{
                \texttt{\tiny{[P1$\langle$GM]}}
                \texttt{Most tegyél fel egy kérdést annak érdekében, hogy új aspektusokat deríts ki az én rácsomról, amely eltérhet a te rácsodtól. Kezdd "KÉRDÉS:" jelzővel, és ne adj hozzá semmi mást.} \\
            }
        }
    }
    & \\ \\

    \theutterance \stepcounter{utterance}  
    & \multicolumn{4}{p{0.6\linewidth}}{
        \cellcolor[rgb]{0.9,0.9,0.9}{
            \makecell[{{p{\linewidth}}}]{
                \texttt{\tiny{[P1$\rangle$GM]}}
                \texttt{KÉRDÉS:} \\
\texttt{A harmadik sorban az X karakterek pontos elhelyezkedése milyen sorrendben van?} \\
            }
        }
    }
    & & \\ \\

    \theutterance \stepcounter{utterance}  
    & & & \multicolumn{2}{p{0.3\linewidth}}{
        \cellcolor[rgb]{0.9,0.9,0.9}{
            \makecell[{{p{\linewidth}}}]{
                \texttt{\tiny{[GM$|$GM]}}
                \texttt{continue} \\
            }
        }
    }
    & & \\ \\

    \theutterance \stepcounter{utterance}  
    & & \multicolumn{4}{p{0.6\linewidth}}{
        \cellcolor[rgb]{0.9,0.9,0.9}{
            \makecell[{{p{\linewidth}}}]{
                \texttt{\tiny{[GM$\rangle$P2]}}
                \texttt{VÁLASZ:} \\
\texttt{Az én második soromban az X karakterek az első, második, negyedik és ötödik pozícióban vannak (X X ▢ X X).} \\
\texttt{KÉRDÉS:} \\
\texttt{A harmadik sorban az X karakterek pontos elhelyezkedése milyen sorrendben van?Kezdd a válaszodat a "VÁLASZ:" jelzővel, és ne adj hozzá semmi mást.} \\
            }
        }
    }
    & \\ \\

    \theutterance \stepcounter{utterance}  
    & & & \multicolumn{4}{p{0.6\linewidth}}{
        \cellcolor[rgb]{0.9,0.9,0.9}{
            \makecell[{{p{\linewidth}}}]{
                \texttt{\tiny{[GM$\langle$P2]}}
                \texttt{VÁLASZ:} \\
\texttt{A harmadik soromban az X karakterek az első, harmadik és ötödik pozícióban vannak (X ▢ X ▢ X).} \\
            }
        }
    }
     \\ \\

    \theutterance \stepcounter{utterance}  
    & & & \multicolumn{2}{p{0.3\linewidth}}{
        \cellcolor[rgb]{0.9,0.9,0.9}{
            \makecell[{{p{\linewidth}}}]{
                \texttt{\tiny{[GM$|$GM]}}
                \texttt{continue} \\
            }
        }
    }
    & & \\ \\

    \theutterance \stepcounter{utterance}  
    & & \multicolumn{4}{p{0.6\linewidth}}{
        \cellcolor[rgb]{0.9,0.9,0.9}{
            \makecell[{{p{\linewidth}}}]{
                \texttt{\tiny{[GM$\rangle$P2]}}
                \texttt{Most tegyél fel egy kérdést annak érdekében, hogy új aspektusokat deríts ki az én rácsomról, amely eltérhet a te rácsodtól. Kezdd "KÉRDÉS:" jelzővel, és ne adj hozzá semmi mást.} \\
            }
        }
    }
    & \\ \\

    \theutterance \stepcounter{utterance}  
    & & & \multicolumn{4}{p{0.6\linewidth}}{
        \cellcolor[rgb]{0.9,0.9,0.9}{
            \makecell[{{p{\linewidth}}}]{
                \texttt{\tiny{[GM$\langle$P2]}}
                \texttt{KÉRDÉS:} \\
\texttt{A negyedik sorban az X karakterek az első és az ötödik pozícióban vannak, vagy vannak benne más X{-}ek is?} \\
            }
        }
    }
     \\ \\

    \theutterance \stepcounter{utterance}  
    & & & \multicolumn{2}{p{0.3\linewidth}}{
        \cellcolor[rgb]{0.9,0.9,0.9}{
            \makecell[{{p{\linewidth}}}]{
                \texttt{\tiny{[GM$|$GM]}}
                \texttt{continue} \\
            }
        }
    }
    & & \\ \\

    \theutterance \stepcounter{utterance}  
    & & \multicolumn{4}{p{0.6\linewidth}}{
        \cellcolor[rgb]{0.9,0.9,0.9}{
            \makecell[{{p{\linewidth}}}]{
                \texttt{\tiny{[P1$\langle$GM]}}
                \texttt{VÁLASZ:} \\
\texttt{A harmadik soromban az X karakterek az első, harmadik és ötödik pozícióban vannak (X ▢ X ▢ X).} \\
\texttt{KÉRDÉS:} \\
\texttt{A negyedik sorban az X karakterek az első és az ötödik pozícióban vannak, vagy vannak benne más X{-}ek is?Kezdd a válaszodat a "VÁLASZ:" jelzővel, és ne adj hozzá semmi mást.} \\
            }
        }
    }
    & \\ \\

    \theutterance \stepcounter{utterance}  
    & \multicolumn{4}{p{0.6\linewidth}}{
        \cellcolor[rgb]{0.9,0.9,0.9}{
            \makecell[{{p{\linewidth}}}]{
                \texttt{\tiny{[P1$\rangle$GM]}}
                \texttt{VÁLASZ:} \\
\texttt{A negyedik soromban csak az első és az ötödik pozícióban vannak az X karakterek (X ▢ ▢ ▢ X).} \\
            }
        }
    }
    & & \\ \\

    \theutterance \stepcounter{utterance}  
    & & & \multicolumn{2}{p{0.3\linewidth}}{
        \cellcolor[rgb]{0.9,0.9,0.9}{
            \makecell[{{p{\linewidth}}}]{
                \texttt{\tiny{[GM$|$GM]}}
                \texttt{continue} \\
            }
        }
    }
    & & \\ \\

    \theutterance \stepcounter{utterance}  
    & & \multicolumn{4}{p{0.6\linewidth}}{
        \cellcolor[rgb]{0.9,0.9,0.9}{
            \makecell[{{p{\linewidth}}}]{
                \texttt{\tiny{[P1$\langle$GM]}}
                \texttt{Most tegyél fel egy kérdést annak érdekében, hogy új aspektusokat deríts ki az én rácsomról, amely eltérhet a te rácsodtól. Kezdd "KÉRDÉS:" jelzővel, és ne adj hozzá semmi mást.} \\
            }
        }
    }
    & \\ \\

    \theutterance \stepcounter{utterance}  
    & \multicolumn{4}{p{0.6\linewidth}}{
        \cellcolor[rgb]{0.9,0.9,0.9}{
            \makecell[{{p{\linewidth}}}]{
                \texttt{\tiny{[P1$\rangle$GM]}}
                \texttt{KÉRDÉS:} \\
\texttt{Az ötödik sorban az X karakterek az első és az ötödik pozícióban vannak, vagy vannak benne más X{-}ek is?} \\
            }
        }
    }
    & & \\ \\

    \theutterance \stepcounter{utterance}  
    & & & \multicolumn{2}{p{0.3\linewidth}}{
        \cellcolor[rgb]{0.9,0.9,0.9}{
            \makecell[{{p{\linewidth}}}]{
                \texttt{\tiny{[GM$|$GM]}}
                \texttt{continue} \\
            }
        }
    }
    & & \\ \\

    \theutterance \stepcounter{utterance}  
    & & \multicolumn{4}{p{0.6\linewidth}}{
        \cellcolor[rgb]{0.9,0.9,0.9}{
            \makecell[{{p{\linewidth}}}]{
                \texttt{\tiny{[GM$\rangle$P2]}}
                \texttt{VÁLASZ:} \\
\texttt{A negyedik soromban csak az első és az ötödik pozícióban vannak az X karakterek (X ▢ ▢ ▢ X).} \\
\texttt{KÉRDÉS:} \\
\texttt{Az ötödik sorban az X karakterek az első és az ötödik pozícióban vannak, vagy vannak benne más X{-}ek is?Kezdd a válaszodat a "VÁLASZ:" jelzővel, és ne adj hozzá semmi mást.} \\
            }
        }
    }
    & \\ \\

    \theutterance \stepcounter{utterance}  
    & & & \multicolumn{4}{p{0.6\linewidth}}{
        \cellcolor[rgb]{0.9,0.9,0.9}{
            \makecell[{{p{\linewidth}}}]{
                \texttt{\tiny{[GM$\langle$P2]}}
                \texttt{VÁLASZ:} \\
\texttt{Az ötödik soromban is csak az első és az ötödik pozícióban vannak az X karakterek (X ▢ ▢ ▢ X).} \\
            }
        }
    }
     \\ \\

    \theutterance \stepcounter{utterance}  
    & & & \multicolumn{2}{p{0.3\linewidth}}{
        \cellcolor[rgb]{0.9,0.9,0.9}{
            \makecell[{{p{\linewidth}}}]{
                \texttt{\tiny{[GM$|$GM]}}
                \texttt{continue} \\
            }
        }
    }
    & & \\ \\

    \theutterance \stepcounter{utterance}  
    & & \multicolumn{4}{p{0.6\linewidth}}{
        \cellcolor[rgb]{0.9,0.9,0.9}{
            \makecell[{{p{\linewidth}}}]{
                \texttt{\tiny{[GM$\rangle$P2]}}
                \texttt{Most tegyél fel egy kérdést annak érdekében, hogy új aspektusokat deríts ki az én rácsomról, amely eltérhet a te rácsodtól. Kezdd "KÉRDÉS:" jelzővel, és ne adj hozzá semmi mást.} \\
            }
        }
    }
    & \\ \\

    \theutterance \stepcounter{utterance}  
    & & & \multicolumn{4}{p{0.6\linewidth}}{
        \cellcolor[rgb]{0.9,0.9,0.9}{
            \makecell[{{p{\linewidth}}}]{
                \texttt{\tiny{[GM$\langle$P2]}}
                \texttt{KÉRDÉS:} \\
\texttt{Az első sorban az X karakterek az első és az ötödik pozícióban vannak, vagy vannak benne más X{-}ek is?} \\
            }
        }
    }
     \\ \\

    \theutterance \stepcounter{utterance}  
    & & & \multicolumn{2}{p{0.3\linewidth}}{
        \cellcolor[rgb]{0.9,0.9,0.9}{
            \makecell[{{p{\linewidth}}}]{
                \texttt{\tiny{[GM$|$GM]}}
                \texttt{continue} \\
            }
        }
    }
    & & \\ \\

    \theutterance \stepcounter{utterance}  
    & & \multicolumn{4}{p{0.6\linewidth}}{
        \cellcolor[rgb]{0.9,0.9,0.9}{
            \makecell[{{p{\linewidth}}}]{
                \texttt{\tiny{[P1$\langle$GM]}}
                \texttt{VÁLASZ:} \\
\texttt{Az ötödik soromban is csak az első és az ötödik pozícióban vannak az X karakterek (X ▢ ▢ ▢ X).} \\
\texttt{KÉRDÉS:} \\
\texttt{Az első sorban az X karakterek az első és az ötödik pozícióban vannak, vagy vannak benne más X{-}ek is?Kezdd a válaszodat a "VÁLASZ:" jelzővel, és ne adj hozzá semmi mást.} \\
            }
        }
    }
    & \\ \\

    \theutterance \stepcounter{utterance}  
    & \multicolumn{4}{p{0.6\linewidth}}{
        \cellcolor[rgb]{0.9,0.9,0.9}{
            \makecell[{{p{\linewidth}}}]{
                \texttt{\tiny{[P1$\rangle$GM]}}
                \texttt{VÁLASZ:} \\
\texttt{Az első soromban csak az első és az ötödik pozícióban vannak az X karakterek (X ▢ ▢ ▢ X).} \\
            }
        }
    }
    & & \\ \\

    \theutterance \stepcounter{utterance}  
    & & & \multicolumn{2}{p{0.3\linewidth}}{
        \cellcolor[rgb]{0.9,0.9,0.9}{
            \makecell[{{p{\linewidth}}}]{
                \texttt{\tiny{[GM$|$GM]}}
                \texttt{continue} \\
            }
        }
    }
    & & \\ \\

    \theutterance \stepcounter{utterance}  
    & & \multicolumn{4}{p{0.6\linewidth}}{
        \cellcolor[rgb]{0.9,0.9,0.9}{
            \makecell[{{p{\linewidth}}}]{
                \texttt{\tiny{[P1$\langle$GM]}}
                \texttt{Most tegyél fel egy kérdést annak érdekében, hogy új aspektusokat deríts ki az én rácsomról, amely eltérhet a te rácsodtól. Kezdd "KÉRDÉS:" jelzővel, és ne adj hozzá semmi mást.} \\
            }
        }
    }
    & \\ \\

    \theutterance \stepcounter{utterance}  
    & \multicolumn{4}{p{0.6\linewidth}}{
        \cellcolor[rgb]{0.9,0.9,0.9}{
            \makecell[{{p{\linewidth}}}]{
                \texttt{\tiny{[P1$\rangle$GM]}}
                \texttt{KÉRDÉS:} \\
\texttt{A második sorodban az X karakterek az első, második, negyedik és ötödik pozícióban vannak, de van benne X a harmadik pozícióban is?} \\
            }
        }
    }
    & & \\ \\

    \theutterance \stepcounter{utterance}  
    & & & \multicolumn{2}{p{0.3\linewidth}}{
        \cellcolor[rgb]{0.9,0.9,0.9}{
            \makecell[{{p{\linewidth}}}]{
                \texttt{\tiny{[GM$|$GM]}}
                \texttt{continue} \\
            }
        }
    }
    & & \\ \\

    \theutterance \stepcounter{utterance}  
    & & \multicolumn{4}{p{0.6\linewidth}}{
        \cellcolor[rgb]{0.9,0.9,0.9}{
            \makecell[{{p{\linewidth}}}]{
                \texttt{\tiny{[GM$\rangle$P2]}}
                \texttt{VÁLASZ:} \\
\texttt{Az első soromban csak az első és az ötödik pozícióban vannak az X karakterek (X ▢ ▢ ▢ X).} \\
\texttt{KÉRDÉS:} \\
\texttt{A második sorodban az X karakterek az első, második, negyedik és ötödik pozícióban vannak, de van benne X a harmadik pozícióban is?Kezdd a válaszodat a "VÁLASZ:" jelzővel, és ne adj hozzá semmi mást.} \\
            }
        }
    }
    & \\ \\

    \theutterance \stepcounter{utterance}  
    & & & \multicolumn{4}{p{0.6\linewidth}}{
        \cellcolor[rgb]{0.9,0.9,0.9}{
            \makecell[{{p{\linewidth}}}]{
                \texttt{\tiny{[GM$\langle$P2]}}
                \texttt{VÁLASZ:} \\
\texttt{A második soromban a harmadik pozícióban nincs X (X ▢ ▢ X X).} \\
            }
        }
    }
     \\ \\

    \theutterance \stepcounter{utterance}  
    & & & \multicolumn{2}{p{0.3\linewidth}}{
        \cellcolor[rgb]{0.9,0.9,0.9}{
            \makecell[{{p{\linewidth}}}]{
                \texttt{\tiny{[GM$|$GM]}}
                \texttt{continue} \\
            }
        }
    }
    & & \\ \\

    \theutterance \stepcounter{utterance}  
    & & \multicolumn{4}{p{0.6\linewidth}}{
        \cellcolor[rgb]{0.9,0.9,0.9}{
            \makecell[{{p{\linewidth}}}]{
                \texttt{\tiny{[GM$\rangle$P2]}}
                \texttt{Most hozz döntést. Mit gondolsz: a te rácsod és a másik rács ugyanazt írja le? Írd azt, hogy "DÖNTÉS: azonos rács", ha szerinted ugyanarról a rácsról van szó, vagy azt, hogy "DÖNTÉS: eltérő rács", ha szerinted a rácsok különbözőek. Ne adj hozzá semmi mást.} \\
            }
        }
    }
    & \\ \\

    \theutterance \stepcounter{utterance}  
    & & & \multicolumn{4}{p{0.6\linewidth}}{
        \cellcolor[rgb]{0.9,0.9,0.9}{
            \makecell[{{p{\linewidth}}}]{
                \texttt{\tiny{[GM$\langle$P2]}}
                \texttt{DÖNTÉS: eltérő rács} \\
            }
        }
    }
     \\ \\

    \theutterance \stepcounter{utterance}  
    & & & \multicolumn{2}{p{0.3\linewidth}}{
        \cellcolor[rgb]{0.9,0.9,0.9}{
            \makecell[{{p{\linewidth}}}]{
                \texttt{\tiny{[GM$|$GM]}}
                \texttt{continue} \\
            }
        }
    }
    & & \\ \\

    \theutterance \stepcounter{utterance}  
    & & & \multicolumn{2}{p{0.3\linewidth}}{
        \cellcolor[rgb]{0.9,0.9,0.9}{
            \makecell[{{p{\linewidth}}}]{
                \texttt{\tiny{[GM$|$GM]}}
                \texttt{success} \\
            }
        }
    }
    & & \\ \\

    \theutterance \stepcounter{utterance}  
    & & \multicolumn{4}{p{0.6\linewidth}}{
        \cellcolor[rgb]{0.9,0.9,0.9}{
            \makecell[{{p{\linewidth}}}]{
                \texttt{\tiny{[P1$\langle$GM]}}
                \texttt{VÁLASZ:} \\
\texttt{A második soromban a harmadik pozícióban nincs X (X ▢ ▢ X X).} \\
\texttt{Most hozz döntést. Mit gondolsz: a te rácsod és a másik rács ugyanazt írja le? Írd azt, hogy "DÖNTÉS: azonos rács", ha szerinted ugyanarról a rácsról van szó, vagy azt, hogy "DÖNTÉS: eltérő rács", ha szerinted a rácsok különbözőek. Ne adj hozzá semmi mást.} \\
            }
        }
    }
    & \\ \\

    \theutterance \stepcounter{utterance}  
    & \multicolumn{4}{p{0.6\linewidth}}{
        \cellcolor[rgb]{0.9,0.9,0.9}{
            \makecell[{{p{\linewidth}}}]{
                \texttt{\tiny{[P1$\rangle$GM]}}
                \texttt{DÖNTÉS: eltérő rács} \\
            }
        }
    }
    & & \\ \\

    \theutterance \stepcounter{utterance}  
    & & & \multicolumn{2}{p{0.3\linewidth}}{
        \cellcolor[rgb]{0.9,0.9,0.9}{
            \makecell[{{p{\linewidth}}}]{
                \texttt{\tiny{[GM$|$GM]}}
                \texttt{continue} \\
            }
        }
    }
    & & \\ \\

    \theutterance \stepcounter{utterance}  
    & & & \multicolumn{2}{p{0.3\linewidth}}{
        \cellcolor[rgb]{0.9,0.9,0.9}{
            \makecell[{{p{\linewidth}}}]{
                \texttt{\tiny{[GM$|$GM]}}
                \texttt{success} \\
            }
        }
    }
    & & \\ \\

\end{supertabular}
}

\end{document}
